将当前向量移动至原点后顺时针旋转 $90^{\circ}$ ，即获取垂直于当前向量的、起点为原点的向量。在计算垂线时非常有用。例如，要想获取点 $a$ 绕点 $o$ 顺时针旋转 $90^{\circ}$ 后的点，可以这样书写代码：`auto ans = o + rotate(o, a);` ；如果是逆时针旋转，那么只需更改符号即可：`auto ans = o - rotate(o, a);` 。